\subsection{Initial and Boundary Conditions}

A complete mathematical formulation requires initial and boundary conditions. Initial conditions define the state of the medium at the beginning of the observation interval, \(t=0\). For the mechanical field, the initial displacement may be prescribed as
\begin{equation}
\mathbf{u}(\mathbf{x},0)
=
\mathbf{u}_0(\mathbf{x}),
\end{equation}
where \(\mathbf{u}_0(\mathbf{x})\) is the initial displacement distribution. The initial velocity may be written as
\begin{equation}
\frac{\partial \mathbf{u}}{\partial t}(\mathbf{x},0)
=
\mathbf{v}_0(\mathbf{x}),
\end{equation}
where \(\mathbf{v}_0(\mathbf{x})\) is the initial velocity field. For the thermal field, the initial temperature distribution is given by
\begin{equation}
T(\mathbf{x},0)
=
T_{\mathrm{init}}(\mathbf{x}).
\end{equation}
These conditions determine the initial thermal and mechanical state from which the transient thermoelastic process begins.

Boundary conditions describe how the selected domain \(\Omega\) interacts with its surroundings through the boundary \(\partial\Omega\). For the mechanical part of the problem, one may prescribe displacement on a part of the boundary \(\Gamma_u\):
\begin{equation}
\mathbf{u}
=
\bar{\mathbf{u}}
\quad
\text{on } \Gamma_u .
\end{equation}
Alternatively, one may prescribe traction on another part of the boundary \(\Gamma_t\):
\begin{equation}
\sigma_{ij}n_j
=
\bar{t}_i
\quad
\text{on } \Gamma_t ,
\end{equation}
where \(\mathbf{n}\) is the outward unit normal vector to the boundary and \(\bar{t}_i\) is the prescribed traction component.

For the thermal part of the problem, temperature may be prescribed on a boundary portion \(\Gamma_T\):
\begin{equation}
T
=
\bar{T}
\quad
\text{on } \Gamma_T .
\end{equation}
Another common thermal condition is prescribed heat flux on \(\Gamma_q\):
\begin{equation}
-k\nabla T \cdot \mathbf{n}
=
\bar{q}
\quad
\text{on } \Gamma_q ,
\end{equation}
where \(\bar{q}\) is the prescribed outward heat flux.

These boundary conditions represent idealized mathematical descriptions of physical constraints. A fixed-displacement boundary may represent a constrained rock surface, while a traction boundary may represent applied stress or free-surface behavior. A prescribed-temperature boundary may represent contact with a heat source or reservoir, while a heat-flux boundary may represent imposed heating, cooling, or insulation. In this thesis, the purpose is not to define a complete field-scale boundary-value problem for a particular geological site, but to identify the types of conditions required for the thermoelastic formulation.